\scanpage{382}
\mspage{298}
% archive sheet number: 381

que $\operatorname{card}\alpha_\pi\deleted{\unsure{>}}=\pi\deleted{\unsure{{}_0|{}'_0}}$,
\deleted{$\operatorname{card}\alpha_0$ est \inserted{\deleted{donc}} le successeur de $\pi'_0$, donc $\leq\pi$.}
L'ens \inserted{ordonné} $I_{\alpha_\pi}$ des
ordinaux $\alpha<\alpha_\pi$ étant g\textsuperscript{r}de devant $\pi'_0$
et de cardinal \deleted{$\unsure{\leq}$} $\pi$, on pourra donc
prendre la limite $\varinjlim_{\alpha\in I_{\alpha_{\unsure{0|\pi}}}}i_\alpha$ dans $I$.
% slanted
\margin{Faux si $\pi$ est un cardinal limite \unsure{!}}

$\circ$ Construction de $(i_\alpha)_{\alpha\in I_{\alpha_0}}$.

\deleted{\unsure{$i_0=i$}} Soit
\[
\begin{aligned}
\alpha_\pi&=\text{plus petit ordinal \inserted{$\alpha$} tel que}\\
&\phantom{=}\operatorname{card}\alpha=\pi
\end{aligned}\tag*{(1)}
\]
donc c'est l'ordinal caractérisé par
\[
\left\{\begin{aligned}
&\operatorname{card}\alpha_\pi=\pi\\
&\alpha<\alpha_\pi\Longrightarrow\operatorname{card}\alpha<\pi,
\end{aligned}\right.\tag*{\unsure{(1)|(2)}}
\]
ou encore par le fait que l'ens ord.
$I_{\alpha_\pi}$ des ordinaux $\alpha<\alpha_\pi$ est \uline{parfaite\textsuperscript{mt}
bien ordonné}. \deleted{\unsure{Donc}} Ainsi, $I_{\alpha_\pi}$ \uline{est
g\textsuperscript{r}d devant $\pi'$ pour tt cardinal} $\pi'<\pi$.
% slanted
\margin{Faux si $\pi$ est cardinal limite \unsure{!}}

Dans les conditions b) \inserted{c)} de 4.13.10, il \unsure{suffit|suffirait}
de supposer $I$ stable par $\varinjlim$ de
type $I_{\alpha_\pi}$, et les $J_{\unsure{r|\alpha}}$ $(r\in A)$
\unsure{stable|stables} dans $I$ par ce type de limites.

\deleted{\unsure{Co}} \uline{Lemme} 4.13.10.1. Soit $\pi$ un cardinal infini,
$A$ un ensemble \inserted{non vide} tel que $\operatorname{card}A\leq\pi$. Alors
